\chapter*{Introduction}

The RSA cryptosystem, named after its inventors Ron Rivest, Adi Shamir, and Leonard Adleman, was first published in 1978 and remains, as of 2026, one of the most widely deployed asymmetric encryption schemes.

Unlike symmetric-key cryptography, where two parties must share a secret key before communicating, RSA allows anyone to encrypt a message using a publicly available key. Only the intended recipient who holds the corresponding private key can decrypt it. The security of RSA relies on the presumed difficulty of factoring the product of two large primes a problem for which no efficient classical algorithm is known, The same mathematical concept (modular exponentiation) also enables digital signatures, providing authentication (proof of origin) and non-repudiation (undeniable authorship).


More than four decades after its invention, RSA is still everywhere. It secures TLS certificates that protect web browsing, authenticates SSH connections, signs software updates, and secures DNSSEC. No single cryptosystem has been more impactful and had a longer time in public-key cryptography.